\begin{mdframed}
    \textbf{La extensión máxima para esta sección es de 4 páginas.} Si utiliza menos páginas y explica bien sus resultados, es preferible.
\end{mdframed}

Presente sus resultados mediante \textbf{tablas y gráficos}. Cada tabla o gráfico debe estar referenciado en el texto, con una breve descripción de lo que se observa y por qué es relevante.

\textbf{Es necesario automatizar la generación de las gráficas}, ya que es imprescindible que se pueda confirmar que las visualizaciones presentadas son producto de los datos generados por sus algoritmos.

Escoja las gráficas más representativas y que muestren de manera clara los resultados obtenidos. Esta elección es parte de lo que se evaluará en la sección de presentación de resultados.



\begin{mdframed}
    Recuerde que es imprescindible que se pueda replicar la generación de las gráficas, por lo que usted debe referir a las figuras generadas en `code/matrix\_multiplication/data/plots/` y `code/sorting/data/plots/`.
\end{mdframed}



% **Algoritmos de Ordenamiento**

% * **Quicksort:** Su tiempo en el mejor caso y en el caso promedio es $O(n \log n)$. Sin embargo, en el peor caso (por ejemplo, si el arreglo ya está ordenado y no se elige un pivote aleatorio), su tiempo decae a $O(n^2)$.
% * **sort (`std::sort` en C++):** En las versiones modernas de C++, se garantiza un tiempo de ejecución de $O(n \log n)$ para el mejor caso, caso promedio y peor caso. Esto se logra usando un algoritmo híbrido llamado Introsort.
% * **Mergesort:** Su tiempo de ejecución es siempre constante en $O(n \log n)$ sin importar si es el mejor, el promedio o el peor caso.

% **Algoritmos de Multiplicación de Matrices**

% * **Naive (Ingenuo):** El método clásico de fuerza bruta (usando tres bucles anidados) tiene un tiempo de ejecución de $O(n^3)$.
% * **Strassen:** Este algoritmo reduce la cantidad de multiplicaciones necesarias, logrando un tiempo de ejecución de $O(n^{\log_2 7})$, lo que equivale aproximadamente a $O(n^{2.81})$. Solo es más rápido que el ingenuo en la práctica cuando las matrices son muy grandes.